A gestão de dados é essencial para a continuidade de atividades de pesquisa, extensão e administração em universidades. No entanto, os conjuntos de dados produzidos nessas instituições frequentemente permanecem dispersos em arquivos locais e sistemas sem integração, carentes de documentação, controle de acesso e preservação de histórico, o que dificulta sua localização, interpretação e reutilização. Diante dessa problemática, o presente trabalho descreve o desenvolvimento do DataRural, uma plataforma web institucional criada para centralizar e organizar a publicação e o compartilhamento de datasets na Universidade Federal Rural do Rio de Janeiro (UFRRJ). A solução permite o envio de conjuntos de dados em formato CSV acompanhados de metadados descritivos, o versionamento com preservação de estados anteriores, o controle de visibilidade entre público e privado e a colaboração por meio de grupos com papéis diferenciados. A plataforma também oferece prévia tabular dos arquivos e download individual ou em pacote. O trabalho demonstra a viabilidade de reunir, em um ambiente institucional, recursos de publicação, rastreabilidade e acesso controlado a dados acadêmicos, e sugere caminhos para aprimoramentos futuros, como suporte a novos formatos, disponibilização de API para consumo programático e avaliação de usabilidade com a comunidade universitária.

\vspace{0.5cm}
\noindent\textbf{Palavras-chave:} gestão de dados, repositório institucional, datasets, plataforma web, versionamento, UFRRJ.